% ============================================================================
% Section 5.3: Practical Deployment Recommendations
% ============================================================================
% Purpose: Translate experimental findings into actionable deployment guidance

\subsection{Practical Deployment Recommendations}
\label{sec:practical}

To ensure our system is immediately deployable, we conducted extensive validation studies to derive actionable guidelines for practitioners facing domain uncertainty in production LLM routing.

% ----------------------------------------------------------------------------
% Recommendation 1: Strategy Selection
% ----------------------------------------------------------------------------
\paragraph{Strategy Selection Based on Prior Quality.}
While Corralling provides robust performance across diverse scenarios, practitioners with knowledge of prior quality can optimize their deployment strategy.
Our experiments demonstrate that when priors are \emph{known} to be severely misspecified (e.g., cross-domain deployment with 68.6\%→13.7\% distribution shift), pure Tabula Rasa achieves 16\% better performance than Corralling (49.5 vs 59.2 cumulative regret).
Conversely, when prior quality is \emph{uncertain}, Corralling provides 18.5\% safety improvement over warmup-only deployment (59.2 vs 74.7 regret), justifying its 20\% overhead relative to the optimal single-expert strategy.

\textbf{Deployment guideline:}
Practitioners should validate prior quality on a small held-out sample from the deployment distribution (N=100--200 requests).
If warmup predictions achieve $>$80\% accuracy, warmup-only deployment is recommended; if $<$50\%, pure Tabula Rasa is optimal; otherwise, Corralling provides robust hedging.
Expert weight evolution (§4.4) provides real-time feedback: warmup weight $<$0.2 after 50 requests indicates harmful priors, while $>$0.8 indicates accurate priors.

% ----------------------------------------------------------------------------
% Recommendation 2: Constant Exploration
% ----------------------------------------------------------------------------
\paragraph{Constant Exploration is Essential Under Mismatch.}
Our ablation study (§4.3) demonstrates that constant exploration ($\alpha=2.0$) is critical for robust performance under domain mismatch, outperforming adaptive decay schedules by 48\% (60.6 vs 90.2 regret).
Standard bandit algorithms employ adaptive decay (e.g., $\alpha_t = \alpha_0/\sqrt{t}$) to transition from exploration to exploitation~\cite{auer2002ucb}.
However, under severe prior mismatch, premature exploitation causes the system to commit irreversibly to misspecified beliefs.

\textbf{Configuration guideline:}
We recommend constant $\alpha=2.0$ for both experts when deploying under domain uncertainty.
This maintains sufficient exploration to detect and adapt to distribution shifts while leveraging available context.
Note that this finding is specific to high-mismatch scenarios; in low-mismatch settings with validated priors, adaptive schedules may improve sample efficiency.
Future work should investigate mismatch-aware $\alpha$ schedules that adjust exploration dynamically based on detected prior quality.

% ----------------------------------------------------------------------------
% Recommendation 3: Monitoring
% ----------------------------------------------------------------------------
\paragraph{Real-Time Monitoring Enables Rapid Iteration.}
Our temporal analysis (§4.4) reveals that expert weight adaptation occurs in approximately 16 $\pm$ 14 requests---significantly faster than anticipated (100--200 requests).
This rapid adaptation is explained by the severity of prior mismatch: when priors are catastrophically wrong, the importance-weighted loss immediately down-weights the harmful expert.

\textbf{Monitoring guideline:}
Practitioners should log expert weights at every request and visualize weight evolution.
Key indicators include:
(1) Warmup weight $<$0.2 indicates harmful priors (consider switching to Tabula Rasa);
(2) Balanced weights (0.3--0.7) indicate genuine uncertainty or transition period;
(3) Warmup weight $>$0.8 indicates accurate priors (consider simplifying to warmup-only).
The high variance in weight trajectories (final warmup weight: 0.382 $\pm$ 0.471) underscores the importance of monitoring \emph{individual} deployments rather than relying on population-level predictions.

% ----------------------------------------------------------------------------
% Summary Table Reference
% ----------------------------------------------------------------------------
Table~\ref{tab:strategy_guide} summarizes our deployment recommendations based on prior quality assessment.
These guidelines enable practitioners to select the optimal routing strategy for their specific deployment context, balancing performance, safety, and implementation complexity.
